\section{Related Work}
\label{sec:related}

\paragraph{The CFR family and exact equilibrium solving.}
CFR \cite{zinkevich2007regret} minimizes counterfactual regret independently at every
information set. In two-player zero-sum games the average strategy profile converges
to a Nash equilibrium \cite{nash1950equilibrium}. CFR$^{+}$
\cite{tammelin2014solving} replaces regret matching \cite{hart2000simple} with regret matching$^{+}$
and linear averaging, and was the engine behind the essential solution of heads-up
limit hold'em \cite{tammelin2015solving,bowling2015heads}.
Subsequent variants tune the update rule: discounting \cite{brown2019solving},
dynamic discounting \cite{xu2024dynamic}, predictive/optimistic updates connecting
CFR to Blackwell approachability
\cite{farina2021faster,farina2019optimistic,farina2019stable,xu2024minimizing}, and
refinements of regret matching itself \cite{farina2023regret,meng2026a}.
First-order methods reach equilibria by a different route, smoothing
\cite{hoda2010smoothing,nesterov2005excessive} on the sequence form
\cite{koller1997representations}. Burch et al.\ \cite{burch2019revisiting} analyze
whether the two players are updated \emph{simultaneously} or \emph{alternately},
which changes iterate quality at matched iteration counts. Our work targets the
execution layer shared by the tabular CFR family. The compiled slot representation
supports vanilla CFR, CFR$^{+}$, DCFR \cite{brown2019solving}, and predictive
CFR$^{+}$ \cite{farina2021faster} under both update orders.

\paragraph{Scaling CFR by approximation.}
When the full tree is too large, the community samples it, shrinks it, prunes it, or
learns it. Monte Carlo CFR estimates counterfactual values from sampled trajectories
\cite{lanctot2009monte}, with public-chance and other variance-controlled variants
for solving and evaluation
\cite{johanson2012efficient,burch2012efficient,schmid2019variance,li2026correlated,li2026av}.
Abstraction maps the game onto a smaller proxy
\cite{gilpin2007lossless,waugh2009abstraction,kroer2018unified,li2024rl,li2025efficient,li2026effective,li2026abstraction},
a line of work motivated directly by CFR's per-iteration cost
\cite{johanson2012efficient,brown2015hierarchical}. Regret-based and online pruning
skip low-value branches during the traversal \cite{brown2015regret,li2025efficient}.
Deep CFR and its successors replace tabular regrets with neural approximation
\cite{brown2019deep,DBLP:journals/corr/abs-1901-07621}, and search-based and
language-model agents combine learned models with online play
\cite{brown2020combining,schmid2023student,li2026pokerskill,wang2026solver}. These methods reduce the game or the
work sampled from it. Our compiler addresses the regime in which the full tree is
available and repeatedly executed. Every one of these methods still runs a tabular CFR
sweep on whatever tree remains, so a faster sweep composes with all of them.

\paragraph{Systems and frameworks for game solving.}
OpenSpiel \cite{lanctot2019openspiel} is the field's reference library. Its CFR
implementations favor readability and breadth of game support. LiteEFG
\cite{liu2024liteefg} compiles extensive-form games into a graph representation
executed by a single-threaded C++ backend. We use it as the primary compiled CPU
baseline. PokerRL \cite{steinberger2019pokerrl} vectorizes CFR over hand ranges
\cite{johanson2012efficient,moravvcik2017deepstack}, a HUNL-specific representation.
Its released Libratus endgames are our external poker reference in
\cref{app:libratus}. Accelerated best-response computation
\cite{johanson2011accelerating} speeds up the \emph{evaluation} half of the loop with
game-specific structure. Our exploitability evaluator is a generic analogue that runs
on the same compiled arrays as the solver. Parallel CFR has been explored through
distributed sampling and equilibrium computation
\cite{jackson2013slumbot,brown2015hierarchical}, on CPU clusters
\cite{DBLP:conf/ispa/ZhangLZXW21a}, and in real-time solvers that partition the
traversal across CPU threads under per-decision budgets \cite{li2026real}. The
milestone poker systems used supercomputer-scale CPU parallelism
\cite{bowling2015heads,brown2018superhuman}. Those designs parallelize the
\emph{traversal}. We change the representation the traversal runs on, so the CPU arm
of our implementation needs no CFR-specific threading of its own (\cref{sec:setup}).

Accelerators already drive game \emph{simulation} at scale \cite{koyamada2023pgx}.
GPU CFR has recent precedent
\cite{kim2026parallelizing,DBLP:journals/corr/abs-2508-06559}, and the GPU path of
the fastest such system still runs slower than an optimized CPU solver. We instead compile
the game into depth-stratified dense tensors and replay the whole iteration as one
captured CUDA graph. This compilation step is what makes stock PyTorch
\cite{paszke2019pytorch} and NVIDIA CUDA graphs \cite{guide2020cuda} applicable to
CFR at all.
